\section{A Fourier estimate for small individual influences}
\label{sec:fourier}

We first work on an arbitrary finite coordinate set $J$. Our objective
is to control low-degree Fourier mass when every individual influence
is small relative to the total. We follow the elementary argument
in~\cite[Lemma~2.1]{Graph2026}; its analytic ingredient is a
two-to-four norm estimate of the type underlying the classical
hypercontractive inequalities~\cite{Bonami70,Beckner75}.
The proof below includes the dependence on the bias.

Fix $0<p<1$ and put $\sigma=\sqrt{p(1-p)}$. For $i\in J$ and
$A\subseteq J$, define
\[
 \chi_i(x)=\frac{x_i-p}{\sigma},\qquad
 \chi_A(x)=\prod_{i\in A}\chi_i(x),\qquad \chi_\varnothing=1.
\]
Independence, zero means, and unit second moments show that the
$\chi_A$ form an orthonormal basis. Write
$\widehat h(A)=\E_p[h\chi_A]$.
Throughout this section, norms and expectations use this product law,
and $I_i(h)$ and $I(h)$ denote pivotal and total influence on $J$.

For Boolean $h$, let $\Delta_i h$ be the function on the other
coordinates obtained by subtracting the section at $x_i=0$ from the
section at $x_i=1$. Direct expectation in coordinate $i$ gives
\[
 \widehat h(A\cup\{i\})=
 \sigma\widehat{\Delta_i h}(A)\qquad(A\subseteq J\setminus\{i\}).
\]
Because $\Delta_i h\in\{-1,0,1\}$, Parseval's identity yields
\begin{equation}\label{eq:parseval}
 \sum_{S\subseteq J}|S|\widehat h(S)^2=\sigma^2 I(h),
 \qquad
 \Var_p(h)=\sum_{\varnothing\ne S\subseteq J}\widehat h(S)^2.
\end{equation}
In particular, this identity uses no sign condition on $\Delta_i h$.

\begin{lemma}[Low-degree mass]\label{lem:low}
Let $m\ge2$ be real and let $h:\{0,1\}^J\to\{0,1\}$ satisfy
$I_i(h)\le I(h)/m$ for every $i\in J$. With
$k=\sigma\log m/64$, one has
\begin{equation}\label{eq:low}
 \sum_{1\le |S|\le k}\widehat h(S)^2\le m^{-1/4}I(h).
\end{equation}
The sum is empty when $k<1$.
\end{lemma}

\begin{proof}
We will use a two-to-four norm estimate to bound low-degree projections
of the differences $\Delta_i h$.
Set $\rho=\sigma/4$ and define the linear operator $T$ by
$T\chi_A=\rho^{|A|}\chi_A$. We first prove
\begin{equation}\label{eq:24}
 \|Tg\|_4\le\|g\|_2
\end{equation}
for every real-valued function $g$ on the product space.
For one coordinate, $|\chi_i|\le1/\sigma$ and $\E\chi_i^2=1$
imply $|\E\chi_i^3|\le1/\sigma$ and
$\E\chi_i^4\le1/\sigma^2$. For real $a,b$, expansion and
$2|ab^3|\le a^2b^2+b^4$ give
\begin{align*}
 \E(a+\rho b\chi_i)^4
 &\le a^4+6\rho^2a^2b^2+
       \frac{2\rho^3}{\sigma}(a^2b^2+b^4)
       +\frac{\rho^4}{\sigma^2}b^4\\
 &=a^4+\frac{13\sigma^2}{32}a^2b^2
       +\frac{9\sigma^2}{256}b^4
 \le(a^2+b^2)^2,
\end{align*}
where the last inequality uses $\sigma\le1/2$.

To extend this estimate to the product, induct on $|J|$, starting with
the constant functions when $J$ is empty. Write $g=a+b\chi_i$,
where $a,b$ depend on the remaining coordinates, and use $T$ for the
operator on those coordinates as well. The one-coordinate estimate,
applied pointwise, followed by the $L^2$ triangle inequality gives
\[
 \|Tg\|_4^2\le\|(Ta)^2+(Tb)^2\|_2
 \le\|Ta\|_4^2+\|Tb\|_4^2
 \le\|a\|_2^2+\|b\|_2^2=\|g\|_2^2.
\]
This proves~\eqref{eq:24}.

Let $P_{\le k}$ be the orthogonal projection onto the span of
$\chi_A$ with $|A|\le k$. For every $u$ in this span,
\[
 \|u\|_4\le\|T^{-1}u\|_2\le\rho^{-k}\|u\|_2.
\]
Taking the supremum of $|\E[gu]|$ over unit $L^2$ vectors in this
span, and applying H\"older's inequality, therefore gives
\begin{equation}\label{eq:projection}
 \|P_{\le k}g\|_2\le\rho^{-k}\|g\|_{4/3}.
\end{equation}
These inequalities are valid for real $k\ge0$, with no rounding
required.

Each nonempty Fourier support is counted at least once when summed
over its coordinates. Using the difference identity and enlarging
the resulting degree cutoff from $k-1$ to $k$, we obtain
\begin{align}
 \sum_{1\le|S|\le k}\widehat h(S)^2
 &\le\sigma^2\sum_{i\in J}
       \|P_{\le k}\Delta_i h\|_2^2 \notag\\
 &\le\sigma^2\rho^{-2k}\sum_{i\in J}I_i(h)^{3/2}.
 \label{eq:derivative-mass}
\end{align}
Here each projection acts on the coordinates other than $i$, and
$\|\Delta_i h\|_{4/3}^2=I_i(h)^{3/2}$ because
$|\Delta_i h|$ is an indicator.

The function $\sigma\log(4/\sigma)$ is increasing on $(0,1/2]$
and is at most $2$. Hence $\rho^{-2k}\le m^{1/16}$.
If $I(h)\le m^{1/4}$, then
$I_i(h)^{1/2}\le m^{-3/8}$, so~\eqref{eq:derivative-mass}
is at most $\sigma^2m^{-5/16}I(h)\le m^{-1/4}I(h)$.
If $I(h)>m^{1/4}$, the right side of~\eqref{eq:low} exceeds $1$,
whereas its left side is at most $\Var_p(h)\le1/4$.
\end{proof}
